% GENERATED FILE. Edit canonical lessons, then run npm run generate:books.
% canonical-source-hash: fnv1a64:2188ad11c5bf1839
% canonical-lessons: GU-C289-adalat, GU-C289-vakil, GU-C289-nyay, GU-C289-saja, GU-C289-bhogavvu

\chapter{In Court}
\label{ch:gu-in-court}
\hlchaptermodality{289}

\begin{chapteropening}
\textbf{By the end of this chapter:} \emph{I can say what the court decided and what the sentence was.}

Complete the last lesson of chapter 289: I can say what the court decided and what the sentence was.
\end{chapteropening}

\section[\texorpdfstring{\gu{અદાલત} (adālat)}{અદાલત (adālat)}]{\gu{અદાલત} (adālat) — a court (of law)}

\label{lesson:GU-C289-adalat}



\begin{quote}
Before the new one: say the Gujarati for a criminal, then the Gujarati for a jail.
\end{quote}

\subsection*{You'll want to know: \gu{અદાલત}}
\textbf{\gu{અદાલત}} — \emph{adālat} — \textquotedblleft{}a court (of law)\textquotedblright{}.

\begin{quote}
\gu{આજે} \gu{સવારે} \gu{પોલીસ} \gu{ગુનેગારને} \gu{અદાલતમાં} \gu{લઈ} \gu{ગઈ}.
\end{quote}

\begin{quote}
\emph{āje savāre polīs gunegārne adālatmā̃ laī gaī.}
\end{quote}

\begin{quote}
This morning the police took the criminal to court.
\end{quote}

A feminine noun.

\begin{cousinweb}
\textbf{\gu{અદાલત}} (\emph{adālat}) is usually traced, through Persian, to Arabic \emph{ʿadāla}, \textquotedblleft{}justice, fairness\textquotedblright{}: the place where justice is done. It is feminine: \textbf{\gu{અદાલતમાં}}, in court.
\end{cousinweb}

\begin{grammarlens}[title={What you've built}]
One more word for telling what happened.
\end{grammarlens}

\subsection*{Your turn}
\begin{itemize}
\raggedright
  \item \emph{Say it:} \emph{adālat}
  \item \emph{Say it:} \emph{adālat}, once more
  \item \emph{Say it:} \emph{āje savāre polīs gunegārne adālatmā̃ laī gaī.} — This morning the police took the criminal to court.
  \item \emph{Recall:} say the Gujarati for a criminal, then the Gujarati for a jail, then say \emph{adālat} again
\end{itemize}

\begin{tcolorbox}[breakable,colback=teal!4,colframe=teal!35!black,title={Before you move on}]
What does \gu{અદાલત} mean? (\textquotedblleft{}A court (of law)\textquotedblright{}.) Say it once more.
\end{tcolorbox}
\section[\texorpdfstring{\gu{વકીલ} (vakīl)}{વકીલ (vakīl)}]{\gu{વકીલ} (vakīl) — a lawyer}

\label{lesson:GU-C289-vakil}



\begin{quote}
Before the new one: say the Gujarati for a jail, then the Gujarati for a court.
\end{quote}

\subsection*{You'll want to know: \gu{વકીલ}}
\textbf{\gu{વકીલ}} — \emph{vakīl} — \textquotedblleft{}a lawyer\textquotedblright{}.

\begin{quote}
\gu{વકીલે} \gu{અદાલતમાં} \gu{ઘણા} \gu{પ્રશ્ન} \gu{પૂછ્યા}.
\end{quote}

\begin{quote}
\emph{vakīle adālatmā̃ ghaṇā praśna pūchyā.}
\end{quote}

\begin{quote}
The lawyer asked many questions in court.
\end{quote}

A masculine or feminine noun, by the person meant.

\begin{cousinweb}
\textbf{\gu{વકીલ}} (\emph{vakīl}) is usually traced, through Persian, to Arabic \emph{wakīl}, \textquotedblleft{}one who acts for another\textquotedblright{}: the person who speaks for you in court.
\end{cousinweb}

\begin{grammarlens}[title={What you've built}]
One more word for telling what happened.
\end{grammarlens}

\subsection*{Your turn}
\begin{itemize}
\raggedright
  \item \emph{Say it:} \emph{vakīl}
  \item \emph{Say it:} \emph{vakīl}, once more
  \item \emph{Say it:} \emph{vakīle adālatmā̃ ghaṇā praśna pūchyā.} — The lawyer asked many questions in court.
  \item \emph{Recall:} say the Gujarati for a jail, then the Gujarati for a court, then say \emph{vakīl} again
\end{itemize}

\begin{tcolorbox}[breakable,colback=teal!4,colframe=teal!35!black,title={Before you move on}]
What does \gu{વકીલ} mean? (\textquotedblleft{}A lawyer\textquotedblright{}.) Say it once more.
\end{tcolorbox}
\section[\texorpdfstring{\gu{ન્યાય} (nyāy)}{ન્યાય (nyāy)}]{\gu{ન્યાય} (nyāy) — justice}

\label{lesson:GU-C289-nyay}



\begin{quote}
Before the new one: say the Gujarati for a court, then the Gujarati for a lawyer.
\end{quote}

\subsection*{You'll want to know: \gu{ન્યાય}}
\textbf{\gu{ન્યાય}} — \emph{nyāy} — \textquotedblleft{}justice\textquotedblright{}.

\begin{quote}
\gu{અદાલતે} \gu{ગરીબ} \gu{માણસને} \gu{ન્યાય} \gu{આપ્યો}.
\end{quote}

\begin{quote}
\emph{adālte garīb māṇasne nyāy āpyo.}
\end{quote}

\begin{quote}
The court gave the poor man justice.
\end{quote}

A masculine noun.

\begin{cousinweb}
\textbf{\gu{ન્યાય}} (\emph{nyāy}) is a tatsama, Sanskrit \emph{nyāya}, \textquotedblleft{}rule, manner, method\textquotedblright{} (CDIAL 7614); in Gujarati it is above all what is fair and right, the thing a court gives.
\end{cousinweb}

\begin{grammarlens}[title={What you've built}]
One more word for telling what happened.
\end{grammarlens}

\subsection*{Your turn}
\begin{itemize}
\raggedright
  \item \emph{Say it:} \emph{nyāy}
  \item \emph{Say it:} \emph{nyāy}, once more
  \item \emph{Say it:} \emph{adālte garīb māṇasne nyāy āpyo.} — The court gave the poor man justice.
  \item \emph{Recall:} say the Gujarati for a court, then the Gujarati for a lawyer, then say \emph{nyāy} again
\end{itemize}

\begin{tcolorbox}[breakable,colback=teal!4,colframe=teal!35!black,title={Before you move on}]
What does \gu{ન્યાય} mean? (\textquotedblleft{}Justice\textquotedblright{}.) Say it once more.
\end{tcolorbox}
\section[\texorpdfstring{\gu{સજા} (sajā)}{સજા (sajā)}]{\gu{સજા} (sajā) — a punishment}

\label{lesson:GU-C289-saja}



\begin{quote}
Before the new one: say the Gujarati for a lawyer, then the Gujarati for justice.
\end{quote}

\subsection*{You'll want to know: \gu{સજા}}
\textbf{\gu{સજા}} — \emph{sajā} — \textquotedblleft{}a punishment\textquotedblright{}.

\begin{quote}
\gu{જો} \gu{અદાલત} \gu{તેને} \gu{ગુનેગાર} \gu{માનશે} \gu{તો} \gu{તેને} \gu{સજા} \gu{થશે}.
\end{quote}

\begin{quote}
\emph{jo adālat tene gunegār mānśe to tene sajā thaśe.}
\end{quote}

\begin{quote}
If the court finds him guilty, he will be punished.
\end{quote}

A feminine noun.

\begin{cousinweb}
\textbf{\gu{સજા}} (\emph{sajā}) is usually traced to Persian \emph{sazā}, \textquotedblleft{}what is fitting, what one deserves\textquotedblright{}. It is feminine; a person \textbf{\gu{સજા} \gu{થવી}}, gets punished, and a court \textbf{\gu{સજા} \gu{આપે} \gu{છે}}, passes sentence.
\end{cousinweb}

\begin{grammarlens}[title={What you've built}]
One more word for telling what happened.
\end{grammarlens}

\subsection*{Your turn}
\begin{itemize}
\raggedright
  \item \emph{Say it:} \emph{sajā}
  \item \emph{Say it:} \emph{sajā}, once more
  \item \emph{Say it:} \emph{jo adālat tene gunegār mānśe to tene sajā thaśe.} — If the court finds him guilty, he will be punished.
  \item \emph{Recall:} say the Gujarati for a lawyer, then the Gujarati for justice, then say \emph{sajā} again
\end{itemize}

\begin{tcolorbox}[breakable,colback=teal!4,colframe=teal!35!black,title={Before you move on}]
What does \gu{સજા} mean? (\textquotedblleft{}A punishment\textquotedblright{}.) Say it once more.
\end{tcolorbox}
\section[\texorpdfstring{\gu{ભોગવવું} (bhogavvũ)}{ભોગવવું (bhogavvũ)}]{\gu{ભોગવવું} (bhogavvũ) — to undergo; to serve (a sentence)}

\label{lesson:GU-C289-bhogavvu}



\begin{quote}
Before the new one: say the Gujarati for justice, then the Gujarati for a punishment.
\end{quote}

\subsection*{You'll want to know: \gu{ભોગવવું}}
\textbf{\gu{ભોગવવું}} — \emph{bhogavvũ} — \textquotedblleft{}to undergo; to serve (a sentence)\textquotedblright{}.

\begin{quote}
\gu{ગુનેગારે} \gu{બે} \gu{વરસની} \gu{જેલ} \gu{ભોગવી}.
\end{quote}

\begin{quote}
\emph{gunegāre be varasnī jel bhogvī.}
\end{quote}

\begin{quote}
The criminal served two years in jail.
\end{quote}

Past, with \gu{તેણે} (he or she …): \gu{તેણે} \gu{સજા} \gu{ભોગવી}, \emph{teṇe sajā bhogvī}.

\begin{cousinweb}
\textbf{\gu{ભોગવવું}} (\emph{bhogavvũ}) is built on Sanskrit \emph{bhoga}, \textquotedblleft{}enjoyment, use\textquotedblright{} (Monier-Williams): to have something to the full. In Gujarati it is also used of what one has to go through, good or bad: a prisoner \textbf{\gu{સજા} \gu{ભોગવે} \gu{છે}}, serves his sentence.
\end{cousinweb}

\begin{grammarlens}[title={What you've built}]
One more word for telling what happened.
\end{grammarlens}

\subsection*{Your turn}
\begin{itemize}
\raggedright
  \item \emph{Say it:} \emph{bhogavvũ}
  \item \emph{Say it:} \emph{bhogavvũ}, once more
  \item \emph{Say it:} \emph{gunegāre be varasnī jel bhogvī.} — The criminal served two years in jail.
  \item \emph{Recall:} say the Gujarati for justice, then the Gujarati for a punishment, then say \emph{bhogavvũ} again
\end{itemize}

\begin{tcolorbox}[breakable,colback=teal!4,colframe=teal!35!black,title={Before you move on}]
What does \gu{ભોગવવું} mean? (\textquotedblleft{}To undergo; to serve (a sentence)\textquotedblright{}.) Say it once more.
\end{tcolorbox}
